\chapter{БАЗОВЫЙ ГРАФИЧЕСКИЙ ИНТЕРФЕЙС ПОЛЬЗОВАТЕЛЯ}
\label{cha:gui}

\section{Постановка задачи}

Графический интерфейс пользователя (ГИП) объединяет подсистемы комплекса в единый
рабочий процесс: создание и сохранение проекта, наполнение сцены компонентами,
управление видом, запуск и контроль симуляции. На данном этапе разработан
\emph{базовый} интерфейс, реализующий основной каркас взаимодействия. Работа
выполнена в репозитории \texttt{modeuler-win} (язык C\#, технология WPF).

\section{Архитектура}

Приложение построено по шаблону Model--View--ViewModel (MVVM), что разделяет
представление (XAML-разметка) и логику (модели представления) и упрощает
тестирование. Команды пользователя связываются с логикой через реализацию
\texttt{RelayCommand} (паттерн <<команда>>).

\textbf{Модели представления} (\texttt{ViewModels/}):
\begin{itemize}
  \item \texttt{MainViewModel} --- корневая модель, координирующая остальные;
  \item \texttt{SceneViewModel} --- состояние трёхмерной сцены и выбранных
        объектов;
  \item \texttt{CatalogueViewModel}, \texttt{GroupedCatalogueViewModel} ---
        каталог компонентов (роботов, объектов окружения);
  \item \texttt{SimulationControlViewModel} --- управление прогоном симуляции
        (старт, пауза, останов);
  \item \texttt{OrientationViewModel} --- индикатор и переключение видовых
        ракурсов;
  \item \texttt{LogPaneViewModel} --- панель журнала событий;
  \item \texttt{ToolBarViewModel} и модели его частей (выбор типа манипуляции,
        привязка к сетке, стыковка панелей).
\end{itemize}

\textbf{Представления} (\texttt{Views/}): главное окно \texttt{MainView} и набор
пользовательских элементов управления --- \texttt{Scene} (область трёхмерного
вида), \texttt{PropertyPane} (свойства выбранного объекта),
\texttt{ComponentCataloguePane} (каталог), \texttt{LogPane} (журнал),
\texttt{Toolbar} (панель инструментов), \texttt{SimulationControl} (управление
симуляцией), \texttt{OrientationControl} (навигационный индикатор).

\textbf{Сервисный слой} (\texttt{Infrastructure/Services/}) скрывает за
интерфейсами реализации функций приложения:
\begin{itemize}
  \item \texttt{ISimulationService} --- запуск и контроль симуляции;
  \item \texttt{ISceneService} --- управление содержимым сцены;
  \item \texttt{IComponentService} --- работа с каталогом компонентов;
  \item \texttt{IGridService}, \texttt{IOrientationService} --- вспомогательная
        геометрия (сетка, ориентация вида);
  \item \texttt{IDockerController} --- запуск вычислительного ядра в контейнере
        и обмен данными с ним;
  \item \texttt{ILoggerPaneService} --- журналирование в интерфейсе.
\end{itemize}

Проект сохраняется и загружается через \texttt{ProjectFileService}
(\texttt{Infrastructure/Persistence}). Свойства компонентов представлены
типизированными классами (\texttt{BooleanProperty}, \texttt{DoubleProperty},
\texttt{IntegerProperty}, \texttt{EnumProperty}, \texttt{StringProperty}), что
позволяет автоматически строить панель свойств по описанию компонента.

\section{Сценарий взаимодействия}

Базовый рабочий процесс: пользователь создаёт проект, выбирает компоненты из
каталога и размещает их на сцене, настраивает их свойства в панели свойств,
управляет видом (ракурсы, сетка, типы манипуляции --- выбор, перемещение,
взаимодействие), после чего запускает симуляцию через панель управления;
ход выполнения и сообщения отображаются в журнале.

\screenshot{0.9}{Главное окно базового графического интерфейса}{Снимок экрана:
главное окно \texttt{modeuler-win} --- область трёхмерного вида в центре, каталог
компонентов и панель свойств по бокам, панель инструментов сверху, панель
журнала и управления симуляцией снизу.}{gui-main}

\section{Результаты}

Реализован базовый графический интерфейс пользователя по архитектуре MVVM с
разделением на представления, модели представления и сервисный слой. Созданы
основные панели (сцена, каталог, свойства, журнал, панель инструментов,
управление симуляцией), механизм проектов и типизированная система свойств
компонентов. Интерфейс предоставляет каркас для интеграции подсистем загрузки,
визуализации и симуляции на последующих этапах.
